\section{Spectral motivation, separator response, and source expansion}
\label{sec:final-spectral-history}

The initial hypothesis concerned audible high-frequency characteristics,
especially in vocals. Candidate measurements represented band-energy ratios,
spectral tilt, entropy, flatness, crest, sibilance-related contrasts and bursts,
and spectral residual structure. They are descriptions of spectral
distribution, not direct measurements of perceived naturalness. Neither high
nor low high-frequency energy is a necessary condition for generated music.

An early metadata negative control reached AUC 1.000 and BA 0.967 using
original container-related information. This demonstrated that the initial
collections could be separated for reasons unrelated to the proposed acoustic
mechanism. High exploratory audio scores were therefore not accepted as
evidence of universal authorship detection. Standardized views, source-aware
evaluation and retained codec/low-pass controls became necessary parts of the
study.

\subsection{Separating the vocal hypothesis from separator error}

The work progressed from mixed audio to separated-vocal analysis through stem
pilots, a 100+100 comparison and expanded 500+500 spectral-band experiments.
Neural source separation used RTX 5090 hardware. A separated vocal is an
estimate, not an oracle reference; leakage and reconstruction error can affect
the same high-frequency measurements that motivated the hypothesis.

The MUSDB18-7 oracle comparison processed 144 excerpts, including fifty
official test tracks. It used a 4,096-point Hann STFT, hop 1,024, 44.1-kHz
audio, a stereo-power magnitude representation, least-squares gain alignment
and an activity floor of $-60$ dB. Signed error was estimated-vocal dB minus
oracle-vocal dB. Frequency/error distributions and representative and severe
STFT-error heatmaps were retained. On the fifty test tracks, median SDR was
8.76 dB and median per-track absolute STFT error approximately 3.28 dB.

An external average frequency-dependent response $b(f)$ was then estimated
without AI/Human labels. The correction applied
\[
 |\widehat V_{\mathrm{corr}}(f,t)|
 = |\widehat V(f,t)|\,10^{-b(f)/20},
\]
retaining the estimated phase. This is a multiplicative response correction
in the magnitude domain, equivalent to a dB-domain offset. It is not the
subtraction of time-aligned waveforms from unrelated songs, nor a general
recovery of separator errors.

Held-out MUSDB median absolute-bin error improved from 3.276 to 2.835 dB.
However, on one independent lossless MedleyDB vocal recording, error increased
from 0.296 to 2.144 dB. That single recording cannot estimate a population
failure rate, but it contradicts a claim of uniformly beneficial correction.
The correction is interpreted as a MUSDB-conditioned stress test.

\begin{table}[htbp]
\centering\small
\caption{Historical 500+500 locked comparison before and after external
separator-response correction. These are not Native30 or unseen-generator
results. Cells give BA / AUC.}
\label{tab:separator-correction}
\begin{tabular}{lrr}
\toprule
Representation & Raw & Corrected\\
\midrule
22 scalar metrics & 0.740 / 0.785 & 0.755 / 0.794\\
0.3--20 kHz band vector & 0.835 / 0.879 & 0.840 / 0.876\\
Band vector excluding 5--10 kHz & 0.845 / 0.922 & 0.855 / 0.922\\
\bottomrule
\end{tabular}
\end{table}

The shared average response did not remove most observed discrimination
(Table~\ref{tab:separator-correction}). This does not exclude class-dependent
separator errors, genre, mastering or codec explanations. A single air-ratio
descriptor had AUC 0.637 in the corrected expanded set, compared with 0.744
in the earlier pilot; this descriptive cross-version contrast cautions against
extrapolating small-sample effects.

\subsection{Generator-specific spectral differences}

The open-model generation work retained 500-prompt experiments for HeartMuLa
and ACE-Step, including prompt and generation provenance. Muse supplied text
prompts; its Suno waveforms were not used as Human recordings. Historical
corrected locked band results were AUC/BA 0.922/0.855 for Suno, 1.000/1.000
for HeartMuLa and 0.988/0.970 for ACE-Step. Each was a generator-specific
task with its own fitted model, not one detector evaluated on all unseen
generators.

Corrected AI-minus-Human mean differences in the 10--16-kHz region were
$+4.322$ dB for Suno, $+14.104$ dB for HeartMuLa and $+8.465$ dB for
ACE-Step. Such heterogeneity motivates generator-held-out evaluation rather
than a single presumed high-frequency AI template. The label-specific
spectral distributions also remain vulnerable to production and collection
confounds.

\subsection{Expanded four-family development ablations}

The expanded S/D/R/P study evaluated all fifteen nonempty subsets at caps
25, 50, 100, 200 and all global groups per source. The per-record spectral
table already covered 7,100 development inputs before the later materialized
expansion; it was not restricted to the original 1,500 AI recordings. The
external MUSDB correction remained fixed rather than being re-estimated from
classification labels. The primary S8 representation restricted analysis below
7.5 kHz with native-rate eligibility; full-band S16 was a separate sensitivity.

\begin{table}[htbp]
\centering\small
\caption{Historical expanded four-family source-transfer results at the all
cap. $J$ and BA are separate macro estimands. Different duration populations
prevent a causal duration interpretation.}
\label{tab:four-family-historical}
\begin{tabular}{lrrrr}
\toprule
Set & 10s $J$ & 10s BA & 30s $J$ & 30s BA\\
\midrule
S & 0.6191 & 0.5903 & 0.7186 & 0.6611\\
D & 0.5487 & 0.5430 & 0.6698 & 0.6370\\
R & 0.4791 & 0.4662 & 0.6566 & 0.6084\\
P & --- & --- & 0.5030 & 0.5153\\
S+D & 0.6334 & 0.5971 & 0.7332 & 0.6711\\
S+D+R & 0.6272 & 0.5931 & 0.7641 & 0.6986\\
S+D+R+P & --- & --- & 0.7601 & 0.6942\\
\bottomrule
\end{tabular}
\end{table}

P was wholly unavailable at ten seconds; combinations containing it were
ineligible for selection. Thirty-second P duration metrics were available on
only 753 rows, with complete bar-span metrics on 414. Missingness was not
interpreted as a measured lack of structural variation.

At caps 25, 50, 100, 200 and all, the ten-second S+D criterion was
0.6198, 0.6226, 0.6297, 0.6317 and 0.6334, respectively. The corresponding
thirty-second S+D+R values were 0.7516, 0.7591, 0.7613, 0.7614 and 0.7641.
Gains flattened at larger caps,
but changing counts, weights and effective regularization prevents attribution
to data quantity alone. Small differences were not supported by CV intervals.

Locked Human specificity exposed a stronger limitation: the thirty-second
candidate reached 95.15\% on MusicNet (330 recordings), but only 52.33\%
on DEAM (300 recordings), corresponding to a 47.67\% DEAM false-positive
rate. Udio family-held-out AUC was 0.5479 and sensitivity 0.1586. These
failures preclude a reliable general-purpose attribution claim despite
favorable within-source spectral separation.

\subsection{Source and output references}

\begingroup\raggedright
Retained report roots under \path{artifacts/} are
\path{frequency_ablation_20260813/REPORT.md},
\path{musdb18_oracle_vocal_eval_20260901/analysis/REPORT.md},
\path{demucs_bias_corrected_1000_20260901/comparison/REPORT.md},
\path{open_models_spectral_500_20260901/analysis/comparison/REPORT.md}, and
\path{source_diversity_expansion_20260905/RESULTS_SUMMARY.md}.
The expanded full fifteen-by-five tables and their numerical audits are in
the latter experiment's \path{results/verified_presentation/} and
\path{audit/formal_results_consistency.json}. These historical results remain
separate from the finalized Exact10 five-family and matched Equal60 analyses.
\par\endgroup
